\documentclass[12pt]{scrartcl}
\usepackage[hagavi]{azzam}

% Aljabar Trigonometri - Modul (gabungan Materi dan Soal Latihan)
% Identitas, persamaan, teleskoping, dan substitusi trigonometri setingkat OSP SMA.
% 25 soal latihan: 12 Identitas, 5 Persamaan, 5 Teleskoping, 3 Substitusi.
% Contoh berlabel OSK/OSP/SMO/Altizio diambil dari arsip Project-Helvethia,
% dilengkapi detail sumber pada komentar % Source. Sisanya komposisi sendiri.
% _Ujian dan _Solusi tetap terpisah (menyusul).

\title{Aljabar Trigonometri - Modul}
\date{6 November 2026}

\begin{document}
\maketitle

\renewcommand*\contentsname{Daftar Isi}
\tableofcontents

\newpage

\part*{Materi}

Modul ini membahas trigonometri dari sisi aljabar: memanipulasi identitas,
menyelesaikan persamaan, meneleskopkan jumlah dan darab, serta memakai
trigonometri sebagai substitusi untuk soal yang awalnya tidak memuat
trigonometri sama sekali. Levelnya OSP SMA. Prasyaratnya hanya definisi
$\sin$, $\cos$, $\tan$ beserta tanda dan relasinya di tiap kuadran.

\section{Identitas Trigonometri}

\begin{definition}
Untuk sudut $x$ yang membuat penyebutnya tidak nol:
\[ \sec x = \frac{1}{\cos x}, \qquad \csc x = \frac{1}{\sin x},
\qquad \cot x = \frac{\cos x}{\sin x} = \frac{1}{\tan x}. \]
Tiga \istilah{identitas Pythagoras}{Pythagorean identity} yang paling sering dipakai:
\[ \sin^2 x + \cos^2 x = 1, \qquad 1 + \tan^2 x = \sec^2 x,
\qquad 1 + \cot^2 x = \csc^2 x. \]
\end{definition}

\begin{example}
Diketahui $\tan x = 2$ dengan $x$ di kuadran III. Tentukan $\sin x + \cos x$.
\end{example}
\begin{proof}[Solusi]
Dari identitas Pythagoras:
\[ \sec^2 x = 1 + \tan^2 x = 1 + 4 = 5. \]
Karena $x$ di kuadran III, $\cos x < 0$ sehingga
\[ \cos x = -\frac{1}{\sqrt{5}}, \qquad
\sin x = \tan x \cos x = -\frac{2}{\sqrt{5}}. \]
Jadi
\[ \sin x + \cos x = -\frac{3}{\sqrt{5}} = -\frac{3\sqrt{5}}{5}. \]
\end{proof}

\begin{theorem}[Jumlah dan selisih]
\istilah{Rumus jumlah dan selisih}{Sum and difference formulas}:
\begin{align*}
\sin(A \pm B) &= \sin A \cos B \pm \cos A \sin B, \\
\cos(A \pm B) &= \cos A \cos B \mp \sin A \sin B, \\
\tan(A \pm B) &= \frac{\tan A \pm \tan B}{1 \mp \tan A \tan B}.
\end{align*}
\end{theorem}

% Source: OSK SMA 2011 isian singkat (sudut lancip A dan B).
\begin{example}[OSK 2011]
Misalkan $A$ dan $B$ sudut lancip dengan $\tan(A+B) = \frac{1}{2}$ dan
$\tan(A-B) = \frac{1}{3}$. Tentukan besar sudut $A$.
\end{example}
\begin{proof}[Solusi]
Tulis $2A = (A+B) + (A-B)$ lalu pakai rumus jumlah tangen:
\begin{align*}
\tan 2A &= \frac{\tan(A+B) + \tan(A-B)}{1 - \tan(A+B)\tan(A-B)} \\
&= \frac{\frac{1}{2} + \frac{1}{3}}{1 - \frac{1}{2} \cdot \frac{1}{3}} \\
&= \frac{5/6}{5/6} = 1.
\end{align*}
Karena $A$ lancip, $2A$ terletak di $(0, \pi)$, sehingga $2A = \frac{\pi}{4}$.
Jadi $A = \frac{\pi}{8} = 22{,}5^{\circ}$.
\end{proof}

% Source: OSP SMA 2009 nomor 10 isian singkat.
\begin{example}[OSP 2009]
Jika $\tan x + \tan y = 25$ dan $\cot x + \cot y = 30$, tentukan $\tan(x+y)$.
\end{example}
\begin{proof}[Solusi]
Karena $\cot$ keduanya terdefinisi, $\tan x$ dan $\tan y$ tidak nol.
Tulis ulang jumlah $\cot$:
\[ \cot x + \cot y = \frac{1}{\tan x} + \frac{1}{\tan y}
= \frac{\tan x + \tan y}{\tan x \tan y}. \]
Substitusi yang diketahui:
\[ 30 = \frac{25}{\tan x \tan y}, \]
sehingga $\tan x \tan y = \frac{5}{6}$. Maka
\[ \tan(x+y) = \frac{\tan x + \tan y}{1 - \tan x \tan y}
= \frac{25}{1 - 5/6} = 150. \]
\end{proof}

\begin{theorem}[Sudut ganda, setengah, dan rangkap tiga]
\istilah{Rumus sudut ganda}{Double-angle formulas}:
\begin{align*}
\sin 2x &= 2\sin x \cos x, \\
\cos 2x &= \cos^2 x - \sin^2 x = 2\cos^2 x - 1 = 1 - 2\sin^2 x, \\
\tan 2x &= \frac{2\tan x}{1 - \tan^2 x}.
\end{align*}
Versi setengah sudutnya:
\[ \sin^2 \frac{x}{2} = \frac{1 - \cos x}{2}, \qquad
\cos^2 \frac{x}{2} = \frac{1 + \cos x}{2}. \]
Versi rangkap tiganya:
\[ \sin 3x = 3\sin x - 4\sin^3 x, \qquad \cos 3x = 4\cos^3 x - 3\cos x. \]
\end{theorem}

% Source: OSK SMA 2005 isian singkat.
% Catatan: kunci di arsip tertulis 7/16, itu salah. Nilai yang benar 7\sqrt{3}/16.
\begin{example}[OSK 2005]
Hitunglah $\sin^{8}75^{\circ} - \cos^{8}75^{\circ}$.
\end{example}
\begin{proof}[Solusi]
Faktorkan sebagai selisih kuadrat dua kali:
\begin{align*}
\sin^8 x - \cos^8 x
&= (\sin^4 x - \cos^4 x)(\sin^4 x + \cos^4 x).
\end{align*}
Faktor pertama:
\[ \sin^4 x - \cos^4 x = (\sin^2 x - \cos^2 x)(\sin^2 x + \cos^2 x) = -\cos 2x. \]
Faktor kedua:
\begin{align*}
\sin^4 x + \cos^4 x &= (\sin^2 x + \cos^2 x)^2 - 2\sin^2 x \cos^2 x \\
&= 1 - \frac{1}{2}\sin^2 2x.
\end{align*}
Untuk $x = 75^{\circ}$ diperoleh $-\cos 150^{\circ} = \frac{\sqrt{3}}{2}$ dan
$1 - \frac{1}{2}\sin^2 150^{\circ} = 1 - \frac{1}{8} = \frac{7}{8}$. Jadi
\[ \sin^{8}75^{\circ} - \cos^{8}75^{\circ}
= \frac{\sqrt{3}}{2} \cdot \frac{7}{8} = \frac{7\sqrt{3}}{16}. \]
\end{proof}

% Source: D. Altizio, Homemade Problems (gaya AIME), diterjemahkan dari bahasa Inggris.
% Arsip: Questions_Bank/Homemade_Problems_David_Altizio.tex di Project-Helvethia.
\begin{example}[D. Altizio, Homemade]
Misalkan $\cos A + \cos B = \frac{5}{4}$ dan $\cos 2A + \cos 2B = \frac{1}{9}$.
Jika $\cos 3A + \cos 3B = \frac{m}{n}$ dalam bentuk paling sederhana,
tentukan $m+n$.
\end{example}
\begin{proof}[Solusi]
Tulis $u = \cos A$ dan $v = \cos B$. Diketahui $u + v = \frac{5}{4}$.
Dari rumus sudut ganda kosinus:
\[ \frac{1}{9} = \cos 2A + \cos 2B = 2u^2 - 1 + 2v^2 - 1, \]
sehingga
\[ u^2 + v^2 = \frac{19}{18}. \]
Maka
\[ uv = \frac{(u+v)^2 - (u^2+v^2)}{2}
= \frac{25/16 - 19/18}{2} = \frac{73}{288}. \]
Selanjutnya
\begin{align*}
u^3 + v^3 &= (u+v)^3 - 3uv(u+v) \\
&= \frac{125}{64} - 3 \cdot \frac{73}{288} \cdot \frac{5}{4} \\
&= \frac{125}{64} - \frac{365}{384} = \frac{385}{384}.
\end{align*}
Terakhir, dengan rumus rangkap tiga:
\begin{align*}
\cos 3A + \cos 3B &= 4(u^3+v^3) - 3(u+v) \\
&= 4 \cdot \frac{385}{384} - \frac{15}{4} \\
&= \frac{385}{96} - \frac{360}{96} = \frac{25}{96}.
\end{align*}
Jadi $m + n = 25 + 96 = 121$.
\end{proof}

\begin{theorem}[Jumlah ke kali dan sebaliknya]
\istilah{Rumus jumlah ke perkalian}{sum-to-product} (selanjutnya disebut sum-to-product):
\begin{align*}
\sin P + \sin Q &= 2\sin\frac{P+Q}{2}\cos\frac{P-Q}{2}, \\
\sin P - \sin Q &= 2\cos\frac{P+Q}{2}\sin\frac{P-Q}{2}, \\
\cos P + \cos Q &= 2\cos\frac{P+Q}{2}\cos\frac{P-Q}{2}, \\
\cos P - \cos Q &= -2\sin\frac{P+Q}{2}\sin\frac{P-Q}{2}.
\end{align*}
\istilah{Rumus perkalian ke jumlah}{product-to-sum} (selanjutnya disebut product-to-sum):
\begin{align*}
\sin A \cos B &= \tfrac12[\sin(A+B) + \sin(A-B)], \\
\cos A \cos B &= \tfrac12[\cos(A+B) + \cos(A-B)], \\
\sin A \sin B &= \tfrac12[\cos(A-B) - \cos(A+B)].
\end{align*}
\end{theorem}

\begin{example}
Hitunglah $\sin 20^{\circ} \sin 40^{\circ} \sin 80^{\circ}$.
\end{example}
\begin{proof}[Solusi]
Ubah darab dua sinus pertama menjadi jumlah:
\[ \sin 20^{\circ} \sin 40^{\circ}
= \tfrac12[\cos 20^{\circ} - \cos 60^{\circ}]
= \tfrac12\cos 20^{\circ} - \tfrac14. \]
Kalikan dengan $\sin 80^{\circ}$:
\[ P = \tfrac12\cos 20^{\circ}\sin 80^{\circ} - \tfrac14\sin 80^{\circ}. \]
Ubah lagi $\cos 20^{\circ}\sin 80^{\circ} = \tfrac12[\sin 100^{\circ} + \sin 60^{\circ}]$
dengan $\sin 100^{\circ} = \sin 80^{\circ}$:
\begin{align*}
P &= \tfrac12\left(\tfrac12\sin 80^{\circ} + \tfrac{\sqrt{3}}{4}\right)
- \tfrac14\sin 80^{\circ} \\
&= \frac{\sqrt{3}}{8}.
\end{align*}
\end{proof}

\begin{fact}[Nilai ekstrem bentuk trigonometri]
Setiap $a\sin x + b\cos x$ dapat ditulis dengan \istilah{sudut bantu}{auxiliary angle}
$\varphi$ sebagai
\[ a\sin x + b\cos x = R\sin(x+\varphi), \]
dengan $R = \sqrt{a^2+b^2}$, $\cos\varphi = a/R$, $\sin\varphi = b/R$.
Akibatnya nilai maksimumnya $R$ dan minimumnya $-R$.
Untuk jumlah yang memuat darab, ketaksamaan $2pq \le p^2 + q^2$ sering
memberi batas atas dan bawah sekaligus.
\end{fact}

% Source: OSK SMA 2009 nomor 14 isian singkat.
% Arsip juga memuat varian (f(y))^2 (OSK SMA 2009 versi 2 nomor 15) dengan jawaban sama.
\begin{example}[OSK 2009]
Misalkan $f(y) = (\sqrt{3}+1)\sin y + (\sqrt{3}-1)\cos y$.
Nilai maksimum $(f(f(y)))^2$ untuk $y$ real adalah \dots
\end{example}
\begin{proof}[Solusi]
Amplitudo $f$ adalah
\begin{align*}
R^2 &= (\sqrt{3}+1)^2 + (\sqrt{3}-1)^2 \\
&= (4+2\sqrt{3}) + (4-2\sqrt{3}) = 8,
\end{align*}
jadi $R = 2\sqrt{2}$ dan daerah hasil $f$ tepat interval $[-2\sqrt{2}, 2\sqrt{2}]$.
Tulis $f(f(y)) = g(t)$ dengan $t = f(y)$ dan $g(t) = a\sin t + b\cos t$.
Nilai terbesar $g$ adalah $R$, dicapai di $t^{\star}$ dengan
\[ \tan t^{\star} = \frac{a}{b} = \frac{\sqrt{3}+1}{\sqrt{3}-1} = 2+\sqrt{3}, \]
yaitu $t^{\star} = \frac{5\pi}{12} \approx 1{,}31$.
Karena $t^{\star}$ termuat di $[-2\sqrt{2}, 2\sqrt{2}]$ dan $f$ kontinu,
ada $y$ dengan $f(y) = t^{\star}$. Jadi maksimumnya $R^2 = 8$.
\end{proof}

% Source: OSP SMA 2009 nomor 15 isian singkat.
\begin{example}[OSP 2009]
Jika $x_1, x_2, \dots, x_{2009}$ bilangan real, nilai terkecil dari
\[ \cos x_1 \sin x_2 + \cos x_2 \sin x_3 + \dots + \cos x_{2009} \sin x_1 \]
adalah \dots
\end{example}
\begin{proof}[Solusi]
Namai jumlah itu $S$. Dari $pq \le \frac{p^2+q^2}{2}$:
\[ S \le \frac{1}{2}\sum_{i=1}^{2009}(\cos^2 x_i + \sin^2 x_{i+1})
= \frac{1}{2}\sum_{i=1}^{2009}(\cos^2 x_i + \sin^2 x_i)
= \frac{2009}{2}, \]
dengan $x_{2010} = x_1$. Terapkan hal yang sama pada $-S$:
\[ -S = \sum_{i=1}^{2009}\cos x_i \sin(-x_{i+1}) \le \frac{2009}{2}, \]
sehingga $S \ge -\frac{2009}{2}$.
Batas ini tercapai, misalnya semua $x_i = \frac{3\pi}{4}$ memberi tiap suku
$\left(-\frac{\sqrt{2}}{2}\right)\left(\frac{\sqrt{2}}{2}\right) = -\frac12$.
Jadi nilai terkecilnya $-\frac{2009}{2}$.
\end{proof}

\section{Persamaan Trigonometri}

\begin{fact}[Bentuk umum solusi dan strategi]
\istilah{Bentuk umum solusi}{general solution} persamaan dasar, dengan $k \in \mathbb{Z}$:
\begin{align*}
\sin x = \sin a &\iff x = a + 2k\pi \text{ atau } x = \pi - a + 2k\pi, \\
\cos x = \cos a &\iff x = a + 2k\pi \text{ atau } x = -a + 2k\pi, \\
\tan x = \tan a &\iff x = a + k\pi.
\end{align*}
Tiga strategi yang paling sering berhasil:
\begin{enumerate}
    \item Samakan bentuk kedua ruas (misalnya semua menjadi sinus), lalu faktorkan.
    \item Kuadratkan untuk menghilangkan akar atau tanda mutlak, lalu wajib cek ulang
    karena penguadratan dapat memunculkan solusi palsu.
    \item Untuk pertidaksamaan, cari batasnya dengan AM-GM atau kuadrat sempurna.
    Tanda kesetaraannya yang menentukan solusi.
\end{enumerate}
\end{fact}

\begin{example}
Tentukan semua bilangan real $x$ dengan $\sin 2x = \sin x$.
\end{example}
\begin{proof}[Solusi]
Langsung dari bentuk umum solusi sinus:
\[ 2x = x + 2k\pi \quad \text{atau} \quad 2x = \pi - x + 2k\pi. \]
Dari ruas pertama $x = 2k\pi$. Dari ruas kedua $3x = \pi + 2k\pi$, yaitu
$x = \frac{\pi}{3} + \frac{2k\pi}{3}$. Semua solusi:
\[ x = 2k\pi \quad \text{atau} \quad x = \frac{\pi}{3} + \frac{2k\pi}{3},
\qquad k \in \mathbb{Z}. \]
\end{proof}

% Source: OSK SMA 2013 isian singkat.
\begin{example}[OSK 2013]
Banyaknya $\alpha$ dengan $0^{\circ} < \alpha < 90^{\circ}$ yang memenuhi
\[ (1+\cos\alpha)(1+\cos 2\alpha)(1+\cos 4\alpha) = \frac18 \]
adalah \dots
\end{example}
\begin{proof}[Solusi]
Tulis $1 + \cos\theta = 2\cos^2\frac{\theta}{2}$:
\[ 8\cos^2\frac{\alpha}{2}\cos^2\alpha\cos^2 2\alpha = \frac18, \]
sehingga
\[ \left|\cos\frac{\alpha}{2}\cos\alpha\cos 2\alpha\right| = \frac18. \]
Karena $0^{\circ} < \alpha < 90^{\circ}$, nilai $\sin\frac{\alpha}{2}$ positif.
Kalikan kedua ruas lalu teleskopkan:
\begin{align*}
\sin\frac{\alpha}{2}\cos\frac{\alpha}{2}\cos\alpha\cos 2\alpha
&= \tfrac12\sin\alpha\cos\alpha\cos 2\alpha \\
&= \tfrac14\sin 2\alpha\cos 2\alpha = \tfrac18\sin 4\alpha.
\end{align*}
Jadi $|\sin 4\alpha| = \sin\frac{\alpha}{2}$, yaitu $\sin 4\alpha = \pm\sin\frac{\alpha}{2}$.
Kasus $\sin 4\alpha = \sin\frac{\alpha}{2}$ memberi $4\alpha = \frac{\alpha}{2} + 360^{\circ}k$
(tidak ada solusi di $(0^{\circ}, 90^{\circ})$) atau
\[ 4\alpha = 180^{\circ} - \frac{\alpha}{2} + 360^{\circ}k, \]
yaitu $\alpha = 40^{\circ} + 80^{\circ}k$, valid untuk $\alpha = 40^{\circ}$.
Kasus $\sin 4\alpha = -\sin\frac{\alpha}{2} = \sin\left(-\frac{\alpha}{2}\right)$
memberi $4\alpha = -\frac{\alpha}{2} + 360^{\circ}k$, yaitu $\alpha = 80^{\circ}k$,
valid untuk $\alpha = 80^{\circ}$; atau
\[ 4\alpha = 180^{\circ} + \frac{\alpha}{2} + 360^{\circ}k, \]
yaitu $\alpha = \frac{360^{\circ}}{7} + \frac{720^{\circ}}{7}k$,
valid untuk $\alpha = \frac{360^{\circ}}{7} \approx 51{,}43^{\circ}$.
Semua langkah di atas dapat dibalik, jadi tepat ada $3$ nilai $\alpha$.
\end{proof}

% Source: OSP SMA 2016 nomor 5 isian singkat.
\begin{example}[OSP 2016]
Jika $0 < x < \frac{\pi}{2}$ dan $4\tan x + 9\cot x \le 12$,
nilai $\sin x$ yang mungkin adalah \dots
\end{example}
\begin{proof}[Solusi]
Tulis $t = \tan x > 0$. Dari AM-GM:
\[ 4t + \frac{9}{t} \ge 2\sqrt{4t \cdot \frac{9}{t}} = 12. \]
Diketahui ruas kiri $\le 12$, jadi harus terjadi kesetaraan:
\[ 4t = \frac{9}{t}, \]
sehingga $t = \frac32$ (yang positif). Karena $x$ lancip:
\[ \sin x = \frac{\tan x}{\sqrt{1+\tan^2 x}}
= \frac{3/2}{\sqrt{13}/2} = \frac{3}{\sqrt{13}} = \frac{3\sqrt{13}}{13}. \]
\end{proof}

\section{Teleskoping Trigonometri}

\begin{theorem}[Teleskoping darab]
Untuk $\sin\frac{x}{2^n} \ne 0$:
\[ \cos\frac{x}{2}\cos\frac{x}{4}\cdots\cos\frac{x}{2^n}
= \frac{\sin x}{2^n\sin\frac{x}{2^n}}. \]
Buktinya: kalikan kedua ruas dengan $\sin\frac{x}{2^n}$ lalu pakai
$2\sin t\cos t = \sin 2t$ berulang kali dari faktor terkecil.
Akibat yang sama guna dari rumus $\tan 2t$:
\[ 1 - \tan^2 t = \frac{2\tan t}{\tan 2t}, \]
sehingga darab bentuk $1 - \tan^2$ juga meneleskop.
\end{theorem}

% Source: OSK SMA 2010 isian singkat.
\begin{example}[OSK 2010]
Nilai $x$ dengan $0 \le x \le \pi$ yang memenuhi
\[ \frac{1}{\sin\frac{x}{2^{2010}}}
= 2^{2010}\sqrt{2}\cos\frac{x}{2}\cos\frac{x}{4}\cdots\cos\frac{x}{2^{2010}} \]
adalah \dots
\end{example}
\begin{proof}[Solusi]
Ruas kiri menuntut $\sin\frac{x}{2^{2010}} \ne 0$, jadi rumus teleskoping boleh dipakai.
Ruas kanan menjadi
\begin{align*}
2^{2010}\sqrt{2} \cdot \frac{\sin x}{2^{2010}\sin\frac{x}{2^{2010}}}
= \frac{\sqrt{2}\sin x}{\sin\frac{x}{2^{2010}}}.
\end{align*}
Persamaannya menyederhanakan menjadi $1 = \sqrt{2}\sin x$, yaitu
$\sin x = \frac{\sqrt{2}}{2}$. Pada $0 \le x \le \pi$:
\[ x = \frac{\pi}{4} \quad \text{atau} \quad x = \frac{3\pi}{4}. \]
Keduanya membuat penyebut $\sin\frac{x}{2^{2010}}$ tidak nol, jadi valid.
\end{proof}

% Source: OSP SMA 2011 nomor 8 isian singkat.
\begin{example}[OSP 2011]
Jika
\[ \left(1-\tan^2\frac{x}{2^{2011}}\right)\left(1-\tan^2\frac{x}{2^{2010}}\right)
\cdots\left(1-\tan^2\frac{x}{2}\right) = 2^{2011}\sqrt{3}\tan\frac{x}{2^{2011}}, \]
maka $\sin 2x$ adalah \dots
\end{example}
\begin{proof}[Solusi]
Tulis tiap faktor sebagai $1 - \tan^2 t = \frac{2\tan t}{\tan 2t}$.
Darabnya meneleskop:
\begin{align*}
\prod_{k=1}^{2011}\left(1-\tan^2\frac{x}{2^k}\right)
&= \prod_{k=1}^{2011}\frac{2\tan\frac{x}{2^k}}{\tan\frac{x}{2^{k-1}}} \\
&= 2^{2011}\frac{\tan\frac{x}{2^{2011}}}{\tan x}.
\end{align*}
Persamaan menjadi
\[ 2^{2011}\frac{\tan\frac{x}{2^{2011}}}{\tan x}
= 2^{2011}\sqrt{3}\tan\frac{x}{2^{2011}}. \]
Jika $\tan\frac{x}{2^{2011}} = 0$, ruas kanan $0$ sedangkan tiap faktor ruas kiri
bernilai $1$, kontradiksi. Jadi boleh dibagi, menghasilkan $\tan x = \frac{1}{\sqrt{3}}$.
Tulis $t = \tan x$:
\[ \sin 2x = \frac{2t}{1+t^2} = \frac{2/\sqrt{3}}{1+1/3} = \frac{\sqrt{3}}{2}. \]
\end{proof}

\begin{theorem}[Teleskoping jumlah]
Untuk $\sin\frac{x}{2} \ne 0$:
\begin{align*}
\sin x + \sin 2x + \dots + \sin nx
&= \frac{\sin\frac{nx}{2}\sin\frac{(n+1)x}{2}}{\sin\frac{x}{2}}, \\
\cos x + \cos 2x + \dots + \cos nx
&= \frac{\sin\frac{nx}{2}\cos\frac{(n+1)x}{2}}{\sin\frac{x}{2}}.
\end{align*}
Buktinya: kalikan jumlahnya dengan $2\sin\frac{x}{2}$, ubah tiap suku dengan
product-to-sum, suku tengahnya saling menghapus.
\end{theorem}

\begin{example}
Selesaikan persamaan $\sin x + \sin 2x + \sin 3x = 0$.
\end{example}
\begin{proof}[Solusi]
Untuk $\sin\frac{x}{2} \ne 0$ (yaitu $x \ne 2k\pi$), rumus teleskoping memberi
\[ \sin x + \sin 2x + \sin 3x
= \frac{\sin\frac{3x}{2}\sin 2x}{\sin\frac{x}{2}}. \]
Persamaan menjadi $\sin\frac{3x}{2}\sin 2x = 0$, sehingga
\[ x = \frac{2k\pi}{3} \quad \text{atau} \quad x = \frac{k\pi}{2},
\qquad k \in \mathbb{Z}. \]
Cek $x = 2k\pi$ secara terpisah: ketiga sukunya nol, jadi memenuhi dan sudah
tercakup dalam $x = \frac{2k\pi}{3}$. Itulah semua solusinya.
\end{proof}

\section{Substitusi Trigonometri}

\begin{fact}[Substitusi bentuk akar]
\istilah{Substitusi trigonometri}{trigonometric substitution} mengganti variabel
agar identitas Pythagoras melenyapkan akarnya. Pilih interval monoton supaya
korespondensinya satu-satu:
\begin{center}
\begin{tabular}{lll}
Bentuk & Substitusi & Interval \\
\hline
$\sqrt{a^2-x^2}$ & $x = a\sin\theta$ atau $a\cos\theta$ & $\theta \in [-\frac{\pi}{2}, \frac{\pi}{2}]$ atau $[0, \pi]$ \\
$\sqrt{a^2+x^2}$ & $x = a\tan\theta$ & $\theta \in (-\frac{\pi}{2}, \frac{\pi}{2})$ \\
$\sqrt{x^2-a^2}$ & $x = a\sec\theta$ & $\theta \in [0, \frac{\pi}{2}) \cup (\frac{\pi}{2}, \pi]$ \\
\end{tabular}
\end{center}
\end{fact}

% Source: SMO 2008 Senior Round 1 pilihan ganda (kunci (B)).
\begin{example}[SMO 2008 Senior]
Jika $f(x) = x^2 + \sqrt{1-x^2}$ dengan $-1 \le x \le 1$,
tentukan daerah hasil $f(x)$.
\end{example}
\begin{proof}[Solusi]
Substitusi $x = \cos\theta$ dengan $\theta \in [0, \pi]$.
Korespondensinya satu-satu dan $\sin\theta \ge 0$, sehingga
\[ \sqrt{1-x^2} = \sqrt{1-\cos^2\theta} = \sin\theta. \]
Fungsinya menjadi
\[ f = \cos^2\theta + \sin\theta = 1 - \sin^2\theta + \sin\theta. \]
Tulis $t = \sin\theta \in [0, 1]$:
\[ g(t) = 1 + t - t^2 = \frac54 - \left(t-\frac12\right)^2. \]
Nilai terbesar $\frac54$ di $t = \frac12$ (tercapai, misalnya $\theta = \frac{\pi}{6}$)
dan nilai terkecil $1$ di ujung interval. Jadi daerah hasilnya $[1, \frac54]$.
\end{proof}

\begin{theorem}[Substitusi syarat jumlah tangen]
Bilangan positif $x, y, z$ memenuhi $x + y + z = xyz$ jika dan hanya jika
ada segitiga lancip $ABC$ dengan $x = \tan A$, $y = \tan B$, $z = \tan C$.
\end{theorem}
\begin{proof}
Arah ($\Leftarrow$): karena $A+B+C = \pi$, berlaku $\tan(A+B) = -\tan C$, yaitu
\[ \frac{\tan A + \tan B}{1 - \tan A \tan B} = -\tan C. \]
(Penyebutnya tidak nol: $\tan A \tan B = 1$ memaksa $C = \frac{\pi}{2}$,
padahal segitiganya lancip.) Kalikan silang:
\[ \tan A + \tan B + \tan C = \tan A \tan B \tan C. \]

Arah ($\Rightarrow$): ambil $A = \arctan x$ dan $B = \arctan y$, keduanya lancip.
Dari $z(xy-1) = x + y > 0$ diperoleh $xy > 1$, sehingga
\[ \tan(A+B) = \frac{x+y}{1-xy} < 0. \]
Jadi $A + B \in (\frac{\pi}{2}, \pi)$ dan $C = \pi - A - B$ lancip dengan
\[ \tan C = -\tan(A+B) = \frac{x+y}{xy-1} = z. \]
\end{proof}

\begin{example}
Diberikan $x, y, z > 0$ dengan $x + y + z = xyz$. Buktikan bahwa
\[ \sqrt{1+x^2} + \sqrt{1+y^2} + \sqrt{1+z^2} \ge 6. \]
\end{example}
\begin{proof}[Solusi]
Substitusi $x = \tan A$, $y = \tan B$, $z = \tan C$ untuk segitiga lancip $ABC$.
Karena $\sec$ positif di sudut lancip:
\[ \sum\sqrt{1+x^2} = \sec A + \sec B + \sec C. \]
Dari AM-GM:
\[ \sec A + \sec B + \sec C \ge \frac{3}{\sqrt[3]{\cos A \cos B \cos C}}. \]
Batasi penyebutnya. Dari product-to-sum:
\[ \cos A \cos B = \frac{\cos(A+B) + \cos(A-B)}{2}
= \frac{-\cos C + \cos(A-B)}{2} \le \frac{1-\cos C}{2}. \]
Karena $t(1-t) \le \frac14$ untuk semua $t$:
\[ \cos A \cos B \cos C \le \frac{\cos C(1-\cos C)}{2} \le \frac18. \]
Akibatnya $\sec A + \sec B + \sec C \ge \frac{3}{\sqrt[3]{1/8}} = 6$.
Kesetaraan saat $x = y = z = \sqrt{3}$.
\end{proof}

\begin{theorem}[Substitusi syarat jumlah kosinus]
Bilangan $x, y, z \in (0, 1)$ memenuhi $x^2 + y^2 + z^2 + 2xyz = 1$ jika dan
hanya jika ada segitiga $ABC$ dengan $x = \cos A$, $y = \cos B$, $z = \cos C$.
\end{theorem}
\begin{proof}
Arah ($\Leftarrow$): dari $C = \pi - A - B$ diperoleh
$\cos C = -\cos(A+B) = -\cos A\cos B + \sin A\sin B$, yaitu
\[ \cos C + \cos A\cos B = \sin A\sin B. \]
Kuadratkan kedua ruas:
\begin{align*}
\cos^2 C + 2\cos A\cos B\cos C + \cos^2 A\cos^2 B
&= \sin^2 A\sin^2 B \\
&= (1-\cos^2 A)(1-\cos^2 B).
\end{align*}
Suku $\cos^2 A\cos^2 B$ di kedua ruas saling menghapus, memberi
$\cos^2 A + \cos^2 B + \cos^2 C + 2\cos A\cos B\cos C = 1$.

Arah ($\Rightarrow$): ambil $A = \arccos x$ dan $B = \arccos y$, keduanya lancip.
Klaim $\cos(A+B) = -z$. Karena
\[ \cos(A+B) = xy - \sqrt{1-x^2}\sqrt{1-y^2}, \]
klaim ini setara dengan $xy + z = \sqrt{1-x^2}\sqrt{1-y^2}$.
Kedua ruas tak negatif, jadi boleh dikuadratkan:
\[ (xy+z)^2 = (1-x^2)(1-y^2) \iff x^2 + y^2 + z^2 + 2xyz = 1, \]
yang tepat asumsi awal. Dari $\cos(A+B) = -z < 0$ diperoleh
$C = \pi - A - B$ lancip dengan $\cos C = z$.
\end{proof}

\begin{example}
Diberikan $x, y, z \in (0, 1)$ dengan $x^2 + y^2 + z^2 + 2xyz = 1$.
Buktikan bahwa $x + y + z \le \frac32$.
\end{example}
\begin{proof}[Solusi]
Substitusi $x = \cos A$, $y = \cos B$, $z = \cos C$ untuk segitiga $ABC$.
Dari sum-to-product:
\[ \cos A + \cos B = 2\cos\frac{A+B}{2}\cos\frac{A-B}{2} \le 2\sin\frac{C}{2}. \]
Tulis $t = \sin\frac{C}{2}$:
\begin{align*}
\cos A + \cos B + \cos C &\le 2t + (1-2t^2) \\
&= \frac32 - 2\left(t-\frac12\right)^2 \le \frac32.
\end{align*}
Kesetaraan saat $A = B = C = 60^{\circ}$, yaitu $x = y = z = \frac12$.
\end{proof}


\part*{Soal Latihan}

\section{Identitas Trigonometri}

\begin{enumerate}

	% Source: OSK SMA 2004 No. 6 (Pilihan Ganda, disajikan sebagai isian)
	% Jawaban: sin 3 < sin 1 < sin 2
	\item (OSK 2004)
	Besar sudut dinyatakan dalam radian. Urutan $\sin 1$, $\sin 2$, $\sin 3$
	dari yang terkecil adalah \dots

	% Source: OSK SMA 2007 No. 10 (Pilihan Ganda, disajikan sebagai isian)
	% Jawaban: sec x - sin x
	\item (OSK 2007)
	Sederhanakan $\frac{\tan^{2}x+\cos^{2}x}{\sin x+\sec x}$
	(untuk $x$ yang membuat penyebutnya tidak nol).

	% Source: SMO 2019 Senior Round 1 (diterjemahkan)
	% Jawaban: 83
	\item (SMO 2019 Senior)
	Diketahui $\tan x = 5$. Tentukan nilai $\frac{6+\sin 2x}{1+\cos 2x}$.

	% Source: SMO 2014 Senior Round 1 (diterjemahkan)
	% Jawaban: 25
	\item (SMO 2014 Senior)
	Hitunglah $50(\cos 39^{\circ}\cos 21^{\circ}
	+ \cos 129^{\circ}\cos 69^{\circ})$.

	% Source: SMO 2013 Open Round 1 (diterjemahkan)
	% Jawaban: 75
	\item (SMO 2013)
	Misalkan $A = \cos^2 10^{\circ} + \cos^2 50^{\circ}
	- \sin 40^{\circ}\sin 80^{\circ}$. Tentukan nilai $100A$.

	% Source: SMO 2013 Open Round 1 (diterjemahkan)
	% Jawaban: 3
	\item (SMO 2013)
	Misalkan $\tan\alpha$ dan $\tan\beta$ adalah dua akar persamaan
	$x^2-3x-3 = 0$. Tentukan nilai
	\[ \left|\sin^2(\alpha+\beta) - 3\sin(\alpha+\beta)\cos(\alpha+\beta)
	- 3\cos^2(\alpha+\beta)\right|. \]

	% Source: OSP SMA 2009 No. 6 (Isian Singkat)
	% Jawaban: 12
	\item (OSP 2009)
	Nilai minimum dari $f(x) = \frac{9x^{2}\sin^{2}x + 4}{x\sin x}$
	untuk $0 < x < \pi$ adalah \dots

	% Source: OSK SMA 2014 (Isian Singkat)
	% Jawaban: 24
	\item (OSK 2014)
	Untuk $0 < x < \pi$, nilai minimum dari
	$\frac{16\sin^2 x + 9}{\sin x}$ adalah \dots

	% Source: SMO 2009 Open Round 1 (diterjemahkan)
	% Jawaban: 2
	\item (SMO 2009)
	Hitunglah $\frac{\sin 80^{\circ}}{\sin 20^{\circ}}
	- \frac{\sqrt{3}}{2\sin 80^{\circ}}$.

	% Source: SMO 2013 Open Round 1 (diterjemahkan)
	% Jawaban: 1
	\item (SMO 2013)
	Tentukan nilai maksimum dari $\frac{1+\cos x}{\sin x+\cos x+2}$
	untuk $x$ bilangan real.

	% Source: OSK SMA 2015 (Isian Singkat)
	% Jawaban: 3/4
	\item (OSK 2015)
	Untuk setiap bilangan real $a$, didefinisikan $f(a)$ sebagai nilai maksimal
	dari $\left|\sin x+\frac{2}{3+\sin x}+a\right|$.
	Nilai minimal dari $f(a)$ adalah \dots

	% Source: OSP SMA 2014 No. 18 (Isian Singkat), diadaptasi.
	% Aslinya menanyakan banyaknya tripel (a,b,c) berurutan dengan kunci 3.
	% Kunci 3 tersebut menghitung yang tak berurutan; yang berurutan ada 15.
	% Di sini soalnya dibuat eksplisit menanyakan himpunan tak berurutan.
	% Jawaban: {2,4,13}, {2,5,8}, {3,3,7}
	\item (OSP 2014)
	Misalkan $0 < \alpha, \beta, \gamma < \frac{\pi}{2}$ dan
	$\alpha+\beta+\gamma = \frac{\pi}{4}$. Tentukan semua himpunan tak berurutan
	$\{a, b, c\}$ bilangan bulat positif sehingga $\tan\alpha = \frac1a$,
	$\tan\beta = \frac1b$, dan $\tan\gamma = \frac1c$.

\end{enumerate}

\section{Persamaan Trigonometri}

\begin{enumerate}

	% Source: SMO 2009 Senior Round 1 (diterjemahkan, rentang x ditambahkan).
	% Soal asli tanpa rentang ("find x"); di sini dibatasi 0 <= x < 360.
	% Jawaban: 95 atau 275
	\item (SMO 2009 Senior)
	Jika $\frac{\cos 100^{\circ}}
	{1-4\sin 25^{\circ}\cos 25^{\circ}\cos 50^{\circ}} = \tan x^{\circ}$
	dengan $0 \le x < 360$, tentukan $x$.

	% Source: SMO 2011 Junior Round 1 (diterjemahkan)
	% Jawaban: 3
	\item (SMO 2011 Junior)
	Barisan $(a_n)$ didefinisikan dengan $a_1 = 3$ dan
	$a_{n+1} = \frac{\sqrt{3}\,a_n-1}{a_n+\sqrt{3}}$ untuk setiap bilangan asli $n$.
	Tentukan $a_{2011}$.

	% Jawaban: x = pi/4 + k pi/2 atau x = +-2pi/3 + 2k pi
	\item
	Tentukan semua bilangan real $x$ dengan $\cos x + \cos 2x + \cos 3x = 0$.

	% Source: D. Altizio, Homemade Problems (diterjemahkan dari bahasa Inggris).
	% Arsip: Questions_Bank/Homemade_Problems_David_Altizio.tex di Project-Helvethia.
	% Jawaban: 231/25
	\item (D. Altizio, Homemade)
	Untuk setiap bilangan real $x$, misalkan $P(x) = 16x^3-21x$.
	Berapakah jumlah semua nilai $\tan^2\theta$ yang mungkin, jika $\theta$
	memenuhi $P(\sin\theta) = P(\cos\theta)$?

	% Jawaban: x = j pi/3, j bilangan bulat
	\item
	Tentukan semua bilangan real $x$ yang membuat semua sukunya terdefinisi dan
	memenuhi $\tan x + \tan 2x + \tan 3x = \tan x\tan 2x\tan 3x$.

\end{enumerate}

\section{Teleskoping Trigonometri}

\begin{enumerate}

	% Jawaban: 1/8
	\item
	Hitunglah $\cos 20^{\circ}\cos 40^{\circ}\cos 80^{\circ}$.

	% Source: SMO 2009 Open Round 1 (diterjemahkan)
	% Jawaban: 12560 (a = 125, b = 60)
	\item (SMO 2009)
	Bentuk $1000\sin 10^{\circ}\cos 20^{\circ}\cos 30^{\circ}\cos 40^{\circ}$
	dapat disederhanakan menjadi $a\sin b^{\circ}$, dengan $a$ dan $b$ bilangan
	bulat positif dan $0 < b < 90$. Tentukan $100a+b$.

	% Jawaban: 1/2
	\item
	Hitunglah $\cos\frac{\pi}{7} - \cos\frac{2\pi}{7} + \cos\frac{3\pi}{7}$.

	% Jawaban: 1/8
	\item
	Hitunglah $\sin\frac{\pi}{14}\sin\frac{3\pi}{14}\sin\frac{5\pi}{14}$.

	% Source: SMO 2024 Senior Round 1 (diterjemahkan)
	% Jawaban: 506
	\item (SMO 2024 Senior)
	Diketahui $\frac{\cos x}{\sin 3x} - \frac{\sin x}{\cos 3x}
	- 2\cdot\frac{\sin 4x}{\cos 6x} = 2024$.
	Tentukan nilai $\frac{\cos 10x}{\sin 12x}$.

\end{enumerate}

\section{Substitusi Trigonometri}

\begin{enumerate}

	% Jawaban: sqrt(2), dicapai di x = 1/sqrt(2)
	\item
	Tentukan nilai maksimum $x + \sqrt{1-x^2}$ untuk $x \in [-1, 1]$.

	% Source: SMO 2024 Open Round 1 (diterjemahkan)
	% Jawaban: 2
	\item (SMO 2024)
	Diketahui $x, y > 0$ dan $x\sqrt{2-y^2} + y\sqrt{2-x^2} = 2$.
	Tentukan nilai $x^2+y^2$.

	% Jawaban: (a) bukti; (b) minimum 2 sqrt(3) di x = y = z = 1/sqrt(3)
	\item
	Diberikan $x, y, z > 0$ dengan $xy+yz+zx = 1$.
	\begin{enumerate}
		\item Buktikan bahwa ada segitiga $ABC$ dengan $x = \tan\frac{A}{2}$,
		$y = \tan\frac{B}{2}$, $z = \tan\frac{C}{2}$.
		\item Buktikan bahwa $\sqrt{1+x^2} + \sqrt{1+y^2} + \sqrt{1+z^2} \ge 2\sqrt{3}$.
	\end{enumerate}

\end{enumerate}

\section{Referensi}
\begin{itemize}
    \item Arsip soal OSK SMA 2005, 2009, 2010, 2011, 2013 dan OSP SMA 2009,
    2011, 2016 (repositori Project-Helvethia).
    \item D. Altizio, \textit{Homemade Problems} (Questions Bank, Project-Helvethia).
    \item SMO 2008 Senior Round 1.
\end{itemize}


\end{document}
